\begin{figure}[tbp]
\centering
\begin{tikzpicture}[x=1cm,y=1cm,
 cluster/.style={draw=blue!55!black,fill=blue!5,ellipse,
 minimum width=1.12cm,minimum height=.82cm,inner sep=5pt},
 contact/.style={draw=blue!65!black,line width=1pt},
 closed/.style={draw=black!45,dashed,line width=.6pt}]
 \node[cluster] (L) at (-5,0) {$L$};
 \node[cluster] (K1) at (-2.8,0) {$K_1$};
 \node[cluster] (K3) at (-.6,0) {$K_3$};
 \node[cluster] (K4) at (1.9,0) {$K_4$};
 \node[cluster] (R) at (4.1,0) {$R$};
 \node[cluster,draw=black!40,fill=black!3] (K2) at (-2.8,-1.55) {$K_2$};
 \draw[contact] (L)--node[above=3pt] {$f_1$}(K1);
 \draw[contact] (K1)--node[above=3pt] {$f_3$}(K3);
 \draw[contact] (K4)--node[above=3pt] {$f_4$}(R);
 \draw[closed] (K3)--node[above=3pt] {$e$}(K4);
 \draw[closed] (K1)--node[right=4pt] {$f_2$}(K2);
 \node[below=5pt] at (L.south) {left base};
 \node[below=5pt] at (R.south) {right base};
 \node[below=5pt,black!65] at (K2.south) {inactive};
\end{tikzpicture}
\caption{The two ancestry chains at the end of interpolation.
Here the terminal path edge $e$ joins $K_3$ to $K_4$.
Only $f_1,f_3,f_4$ are opened. Each terminal cluster joins its own
large base; the two sides remain distinct. The inactive contact
$f_2$ stays closed. Vertices inside each cluster are suppressed.}
\label{fig:interpolation}
\end{figure}
